\label{sec:conc}

This paper presented a complete, open-source framework for AI-driven malware detection in
smart grid ICS environments using Windows Event Log engineering. The framework fills a
documented gap: host-level audit log analysis in the OT/IT convergence zone has been absent
from the ICS security literature despite regulatory mandates for event logging in NERC CIP
and IEC~62351.

Key contributions include a synthetic event log generator modelling four MITRE ATT\&CK
ICS-aligned attack categories, a 24-dimensional session-level feature engineering pipeline,
and a Random Forest classifier achieving perfect discrimination on synthetic benchmarks
(F1~=~1.00, ROC-AUC~=~1.00). The ablation study established that temporal features
(time-delta statistics, sequence entropy) carry the highest discriminative value, consistent
with the burst-timing signatures of reconnaissance and lateral movement.

Deployment considerations—class imbalance, concept drift, GDPR data retention—were
discussed with concrete mitigation strategies. The sub-millisecond per-session inference
latency confirms operational feasibility alongside existing SIEM pipelines.

Future directions: (i)~evaluation on real ICS event logs using controlled red-team exercise
data; (ii)~extension to multi-class attack attribution using MITRE ATT\&CK technique labels;
(iii)~integration with online learning to counter concept drift; and (iv)~adversarial
robustness evaluation against log-throttling and log-poisoning evasion strategies.
